%% CUP Data Template starter — CUP-JNL-DTM document class.
%%
%% This class file (CUP-JNL-DTM.cls) IS bundled here: it is not part of any
%% TeX Live distribution, so a download without it would not compile. Its
%% own header carries no licence grant for the class itself — see this
%% template's README.md for what that means and why it ships anyway.
%%
%% The class picks your bibliography TOOLCHAIN automatically from the
%% `journal=' option below (natbib+bibtex for some journals, biblatex for
%% others) — do not add your own \usepackage{natbib} or {biblatex}, and do
%% not add \addbibresource yourself; see the Bibliography note at the end of
%% this file for what changes if you pick a biblatex journal.
%%
%%%%%%%%%%%%%%%%%%%%%%%%%%%%%%%%%%%%%%%%%%%%%%%%
% Journal codes for the `journal=' class option, and which toolchain each
% one selects:
%   Glasgow Mathematical Journal: gmj      (bibtex + natbib, amsplain style)
%   Forum of Mathematics, Pi:     fmp      (bibtex + natbib, amsplain style)
%   Forum of Mathematics, Sigma:  fms      (bibtex + natbib, amsplain style)
%   Wearable Technologies:        wet      (biblatex, author-year)
%   Data and Policy:              dap      (biblatex, author-year)
%   Data-centric Engineering:     dce      (biblatex, author-year)
%   Environmental Data Science:   eds      (biblatex, author-year)
%   Programmable Materials:       prm      (biblatex, author-year)
%   Journal of Nonlinear Waves:   jnw      (biblatex, author-year)
%   Flow:                         flw      (biblatex, APA style)
%   Judgment and Decision Making: jdm      (biblatex, APA style)
%   Psychometrika:                psy      (biblatex, APA style)
%   Research Synthesis Methods:   rsm      (biblatex, numbered superscript)
%%%%%%%%%%%%%%%%%%%%%%%%%%%%%%%%%%%%%%%%%%%%%%%%
\documentclass[journal=gmj]{CUP-JNL-DTM}

\usepackage{graphicx}
\usepackage{amsmath,amssymb,amsfonts,amsthm}
\usepackage[T1]{fontenc}
\usepackage[colorlinks,allcolors=blue]{hyperref}

\newtheorem{theorem}{Theorem}[section]
\theoremstyle{definition}
\newtheorem{remark}[theorem]{Remark}

\jname{Glasgow Mathematical Journal}
\articletype{Research Article}
\jyear{2026}

\begin{document}

\begin{Frontmatter}

\title[Convergence of a discretised heat kernel on bounded domains]{Convergence of a discretised heat kernel on bounded domains}

\author[1]{Ada Lovelace}
\author[2]{Grace Hopper}

\authormark{Lovelace and Hopper}

\address[1]{\orgdiv{School of Mathematics}, \orgname{University of Somewhere}, \orgaddress{\city{Glasgow}, \postcode{G12 8QQ}, \country{United Kingdom}}}
\address[2]{\orgdiv{Department of Applied Analysis}, \orgname{Institute of Computing}, \orgaddress{\city{Cambridge}, \postcode{CB2 1TN}, \country{United Kingdom}}. \email{g.hopper@example.ac.uk}}

\keywords{heat kernel, finite differences, spectral convergence, bounded domains}

\keywords[MSC Codes]{\codes[Primary]{35K05}; \codes[Secondary]{65M06, 65M12}}

\abstract{We study a finite-difference discretisation of the heat kernel on
a bounded Lipschitz domain and show that the discrete spectrum converges to
the continuous Dirichlet spectrum at a rate governed by mesh regularity. The
argument combines a discrete maximum principle with a compactness estimate
that avoids explicit boundary-layer correctors, giving a proof that extends
to piecewise-smooth boundaries. Numerical experiments on a square and an
L-shaped domain corroborate the predicted first-order rate.}

\end{Frontmatter}

\section{Introduction}

The heat equation on a bounded domain $\Omega \subset \mathbb{R}^d$ with
Dirichlet boundary conditions is one of the most-studied elliptic problems in
applied analysis, yet the convergence rate of finite-difference schemes on
domains with non-smooth boundaries is still not settled in full generality.
Classical results \citep{knuth:1984} establish first-order convergence for
smooth boundaries; for corner domains the rate typically degrades unless the
mesh is graded near the corner \citep{lesk:1977}.

This paper gives a self-contained convergence proof that (i) applies to any
bounded Lipschitz domain, (ii) does not require an explicit boundary-layer
corrector, and (iii) recovers the classical first-order rate when the
boundary is smooth.

\section{Setting and Main Result}

Let $\Omega \subset \mathbb{R}^d$ be a bounded Lipschitz domain and consider
the heat equation
\begin{equation}
  \partial_t u = \Delta u \quad \text{in } \Omega \times (0, T], \qquad
  u = 0 \text{ on } \partial\Omega, \qquad u(\cdot, 0) = u_0.
  \label{eq:heat}
\end{equation}
We discretise \eqref{eq:heat} on a uniform grid of spacing $h$ restricted to
$\Omega$, using the standard five-point Laplacian and forward Euler in time
with step $\tau = c h^2$ for a fixed Courant number $c \le 1/(2d)$.

\begin{theorem}[Discrete spectral convergence]
\label{thm:main}
Let $\lambda_h$ be the smallest eigenvalue of the discrete Dirichlet
Laplacian on $\Omega_h = \Omega \cap h\mathbb{Z}^d$, and let $\lambda_1$ be
the smallest eigenvalue of the continuous Dirichlet Laplacian on $\Omega$.
Then there is a constant $C$, depending only on $\Omega$ and $d$, such that
\begin{equation}
  |\lambda_h - \lambda_1| \le C h,
\end{equation}
for all sufficiently small $h$.
\end{theorem}

\begin{remark}
The constant $C$ in Theorem~\ref{thm:main} depends on the Lipschitz constant
of $\partial\Omega$ but not on any interior regularity of the eigenfunction,
which is what allows the argument to avoid a boundary-layer corrector.
\end{remark}

\begin{proof}[Proof sketch]
The upper bound follows from a discrete Rayleigh quotient argument applied to
a mollified restriction of the continuous first eigenfunction. The lower
bound uses a discrete maximum principle to control the discrepancy between
$\Omega_h$ and $\Omega$ near the boundary, avoided entirely for smooth
boundaries and handled via a covering argument by half-spaces for Lipschitz
corners. The full argument is given in the appendix.
\end{proof}

\section{Numerical Experiments}

We verify Theorem~\ref{thm:main} on two test domains: the unit square, where
the exact spectrum is known in closed form, and an L-shaped domain, which has
a single reentrant corner.

\begin{table}[t]
\tabcolsep=0pt%
\TBL{\caption{Observed convergence rate for $|\lambda_h - \lambda_1|$ as the
mesh spacing $h$ halves. The L-shaped domain's reentrant corner does not
degrade the rate below first order, consistent with
Theorem~\ref{thm:main}.\label{tab:rates}}}
{\begin{fntable}
\begin{tabular*}{\textwidth}{@{\extracolsep{\fill}}lccc@{}}\toprule%
\TCH{Domain} & \TCH{$h$} & \TCH{$|\lambda_h - \lambda_1|$} & \TCH{Observed rate} \\\midrule
Square    & $2^{-4}$ & $1.83\times 10^{-2}$ & --- \\
Square    & $2^{-5}$ & $9.14\times 10^{-3}$ & $1.00$ \\
Square    & $2^{-6}$ & $4.57\times 10^{-3}$ & $1.00$ \\
L-shape   & $2^{-4}$ & $2.41\times 10^{-2}$ & --- \\
L-shape   & $2^{-5}$ & $1.22\times 10^{-2}$ & $0.98$ \\
L-shape   & $2^{-6}$ & $6.15\times 10^{-3}$ & $0.99$ \\
\botrule
\end{tabular*}%
\footnotetext[]{{Note:} rates are computed as $\log_2$ of the ratio of
successive errors.}
\end{fntable}}
\end{table}

Table~\ref{tab:rates} shows the observed rate holding at approximately $1$ on
both domains, matching the bound in Theorem~\ref{thm:main} even at the
reentrant corner of the L-shaped domain.

\section{Conclusion}

We have shown that the discrete Dirichlet spectrum converges to the
continuous spectrum at first order on any bounded Lipschitz domain, with a
proof that does not require an explicit boundary-layer corrector. The bound
is sharp on both test domains, including one with a reentrant corner. A
natural next step is extending the argument to Neumann boundary conditions,
where the discrete maximum principle used here does not directly apply.

\begin{appendix}
\section{Proof of Theorem~\ref{thm:main}}
\label{appendixA}

We give the full argument omitted from Section~2. Fix $h > 0$ small and let
$u_h$ denote the discrete first eigenfunction on $\Omega_h$, normalised so
that $\|u_h\|_{\ell^2(\Omega_h)} = 1$. Extend $u_h$ to $\Omega$ by piecewise
constant interpolation on the grid cells and denote the extension
$\tilde{u}_h$. A standard discrete Rayleigh quotient estimate, together with
the Lipschitz regularity of $\partial\Omega$, gives
\[
  \lambda_1 \le \frac{\int_\Omega |\nabla \tilde{u}_h|^2}{\int_\Omega \tilde{u}_h^2} \le \lambda_h + C_1 h,
\]
which is the upper bound. The lower bound follows from a discrete maximum
principle argument applied to the difference between the discrete and
continuous resolvents, covering the boundary by a finite collection of
half-space charts (possible since $\partial\Omega$ is Lipschitz), giving a
matching constant $C_2$ with $C = \max(C_1, C_2)$.
\end{appendix}

\begin{Backmatter}

\paragraph{Acknowledgments}
We thank the referees for comments that simplified the boundary argument in
the appendix.

\paragraph{Funding Statement}
This research received no specific grant from any funding agency, commercial
or not-for-profit sectors.

\paragraph{Competing Interests}
None.

\paragraph{Data Availability Statement}
The code used to generate Table~\ref{tab:rates} is available on request from
the corresponding author.

\paragraph{Author Contributions}
Conceptualization: A.L. Methodology: A.L., G.H. Formal analysis: A.L.
Numerical experiments: G.H. Writing: A.L., G.H. Both authors approved the
final submitted draft.

%% Bibliography: this journal code (gmj) uses the class's default bibtex +
%% natbib[numbers] + amsplain toolchain, set automatically by `journal=gmj'
%% above — hence citations rendering as numbers, not author-year, even
%% though \citep{} is the natbib command used below. Only \bibliography{}
%% is needed; do not add \bibliographystyle, the class sets amsplain itself.
%%
%% If you switch `journal=' to a biblatex journal (see the table at the top
%% of this file), replace this \bibliography line with
%% \addbibresource{references.bib} moved up to the preamble (before
%% \begin{document}) and \printbibliography here instead.
\bibliography{references}

\end{Backmatter}

\end{document}
